\documentclass[10pt,a4paper]{article}
\usepackage{amsmath,amssymb,amsthm,graphicx}
% Card-local minimal preamble (journal unit sigma-2025-020-non-integrability-of-the-sasano-system-of-ty, P2-ahl16 precedent).
%% Verbatim from the pinned source: the author theorem environments and macro block,
%% so every macro and environment of the counted spans is available.
% Omitted and why: the documentclass line (the journal class sigma is not on the host;
%% the assembler supplies the class) and page-geometry/font blocks (they fight the
%% forced xelatex engine).
\usepackage{amsmath,amssymb,amsthm,mathtools}
\usepackage{xcolor}
\usepackage{url}
\usepackage[hidelinks]{hyperref}

\newcommand{\RS}[1]{\textcolor{red}{#1}}
\numberwithin{equation}{section}
\newtheorem{Theorem}{Theorem}[section]
\newtheorem*{Theorem*}{Theorem}
\newtheorem{Corollary}[Theorem]{Corollary}
\newtheorem{Lemma}[Theorem]{Lemma}
\newtheorem{Proposition}[Theorem]{Proposition}
 { \theoremstyle{definition}
\newtheorem{Definition}[Theorem]{Definition}
\newtheorem{Note}[Theorem]{Note}
\newtheorem{Example}[Theorem]{Example}
\newtheorem{Remark}[Theorem]{Remark} }
\usepackage{cleveref}
\crefname{originalreference}{original reference}{original references}
\Crefname{originalreference}{Original reference}{Original references}
\makeatletter
\newlabel{for}{{3.1}{source}{}{sourcepointer-1}{}}
\newlabel{for@cref}{{[originalreference][0][]3.1}{[1][0][]source}{}{}{}}
\makeatother
\title{One-summability of the two divergent series in every non-singular direction}
\author{Stoyanova, Tsvetana}
\date{doi:10.3842/SIGMA.2025.020; CC BY 4.0}
\begin{document}
\maketitle
\setcounter{section}{1}
\setcounter{Theorem}{17}
\noindent\textit{Editorial note.} All delivered mathematical text is the author's own, copied from the pinned journal source file sources/sigma-2025-020-non-integrability-of-the-sasano-system-of-ty/pinned-original-source.tex (Non-Integrability of the Sasano System of Type D5(1) and Stokes Phenomena); no statement or proof step was rewritten, reordered or removed.\par\medskip
\noindent\textit{Editorial note.} The counted proof is the author's complete argument: All delivered mathematical text is the author's own, copied verbatim from the pinned SIGMA source file pinned-original-source.tex. Counted statement: Lemma environment at source lines 822-855 (label sum), own proof at 857-922. The counted definition environment is labelled sum, which is the author's own label for the Lemma (not a numbered equation).\par\medskip
\noindent\textit{Editorial note.} The label sum names the counted statement, and the proof environment immediately following it in the source is delivered verbatim as its proof.\par\medskip
\noindent\textit{Editorial note.} The counted statement is carried by the author's own Lemma environment in the article source.\par\medskip
\noindent\textit{Editorial note.} Journal-unit assembly: card-local preamble (author theorem environments and macros verbatim), article class, bibliography backend none; any literature citation carried by the delivered span is disclosed in the notes of this card.\par\medskip
\noindent\textit{Editorial note.} Journal numbering seed recomputed from the source itself; printed number 1.18 per census v3.7 is only a cross-check (section 8.8.98).\par\medskip
\noindent\textit{Editorial note.} Revision r002: regenerated from the same spec after the r001 assembly crashed during packaging.\par\medskip
\noindent\textit{Editorial scope.} the counted statement and its own proof environment only\par\medskip
\noindent\textit{Allowed standard prerequisites (editorial).}\begin{itemize}
\item The divergent power series phi-hat and phi-hat introduced by Propos.... Supplies background that the delivered proof uses. The entry is declared as a dependency; it is not re-proved on the card, and the load-bearing steps of the counted statement are carried out inside the delivered span.
\end{itemize}\medskip
\ifdefined\hypertarget\hypertarget{sourcepointer-1}{}\fi
\noindent\textit{Source reference.} Original reference 3.1: sources/sigma-2025-020-non-integrability-of-the-sasano-system-of-ty/pinned-original-source.tex, source lines 651--651. Navigation pointer invoked at line(s) 858: the counted text refers to this environment of the same journal unit, which is not retained in this card, so the reference is declared as a pointer carrying the number the article itself prints for it (the section and the place of the statement inside that section, derived from the shared section-scoped theorem counter of the source); the cited item is an allowed input and none of it is reproduced inside the span.\par\medskip
 \begin{Lemma}\label{sum}
 Under the above notations, we have
 \begin{enumerate} 	\itemsep=0pt
 \item[$1.$]
 	Assume that $\alpha\notin\mathbb{N}$. Then for every direction $\theta \neq \pi$ the function
 	 \[
 	 \varphi_{\theta}(\tau)=2 (\alpha-1)
 	 \int_0^{+\infty {\rm e}^{{\rm i} \theta}}
 	 (1+\zeta)^{\alpha-2} {\rm e}^{-\frac{\zeta}{\tau}} {\rm d} \zeta
 	 \]
 	 defines the $1$-sum of the power series $\hat{\varphi}(\tau)$ in such a direction.
 	

 \item[$2.$]
 Assume that $\alpha\notin\mathbb{Z}_{\leq 0}$. Then for every direction $\theta \neq 0$ the function
 \[
 \phi_{\theta}(\tau)=\int_0^{+\infty {\rm e}^{{\rm i} \theta}}
 (1-\zeta)^{-\alpha} {\rm e}^{-\frac{\zeta}{\tau}} {\rm d} \zeta
 \]
 defines the $1$-sum of the power series $\hat{\phi}(\tau)$ in such a direction.

 \end{enumerate}

 When ${\rm Re} (\alpha) \leq 2$ $($resp.\ ${\rm Re} (\alpha) \geq 0)$ the function $\varphi_{\theta}(\tau)$
 $($resp.\ $\phi_{\theta}(\theta))$ is a holomorphic function in the open unlimited disc
 \[
 \mathcal{D}_{\theta}=\biggl\{ \tau\in\mathbb{C} \mid {\rm Re} \biggl(\frac{{\rm e}^{{\rm i} \theta}}{\tau}\biggr) > 0 \biggr\}
 \]
 for any direction $\theta \neq \pi$ $($resp.\ $\theta\neq 0)$.
 Otherwise, they are holomorphic functions in the open bounded disc
 \[
 \mathcal{D}_{\theta}(1)=\biggl\{ \tau\in\mathbb{C} \mid {\rm Re} \biggl(\frac{{\rm e}^{{\rm i} \theta}}{\tau}\biggr) > 1 \biggr\} .
 \]
 \end{Lemma}

 \begin{proof}
 Consider the divergent power series $\hat{\varphi}(\tau)$ and $\hat{\phi}(\tau)$ defined by Proposition \ref{for}. For the corresponding
 formal Borel transforms, we obtain the power series
 \begin{gather*}
 \hat{\mathcal{B}}_1 \hat{\varphi}(\zeta)
 =
 2 \sum_{n=0}^{\infty} (\alpha-1)_{n+1} \frac{\zeta^n}{n!}=
 2 (\alpha-1) \sum_{n=0}^{\infty} (\alpha-2)_n \frac{\zeta^n}{n!},\qquad
 \hat{\mathcal{B}}_1 \hat{\phi}(\zeta)
 =
 \sum_{n=0}^{\infty} \alpha^{(n)} \frac{\zeta^n}{n!} .
 \end{gather*}
 Both of the series $\hat{\mathcal{B}}_1 \hat{\varphi}(\zeta)$ and $\hat{\mathcal{B}}_1 \hat{\phi}(\zeta)$ are convergent near the origin
 of the Borel $\zeta$-plane with finite radiuses of convergence. Therefore, both of the divergent series $\hat{\varphi}(\tau)$
 and $\hat{\phi}(\tau)$ are Gevrey series of order 1. The functions
 \[
 \varphi(\zeta)=2 (\alpha-1) (1+\zeta)^{\alpha-2},\qquad
 \phi(\zeta)=(1-\zeta)^{-\alpha}
 \]
 present the sums of the power series $\hat{\mathcal{B}}_1 \hat{\varphi}(\zeta)$ and $\hat{\mathcal{B}}_1 \hat{\phi}(\zeta)$, respectively.

 Consider the function $(1+\zeta)^{\alpha-2}$. We have that $\bigl|(1+\zeta)^{\alpha-2}\bigr|=A |1+\zeta|^{{\rm Re}(\alpha) -2}$, where
 $A={\rm e}^{-{\rm Im} (\alpha) \arg (1+\zeta)}$. Let $\theta=\arg (\zeta)$. If ${\rm Re} (\alpha) -2 \leq 0$, then
 \[
 \frac{A}{|1+\zeta|^{2 - {\rm Re}(\alpha)}} \leq A
 \]
 when $\cos \theta \geq 0$, while
 \[
 \frac{A}{|1+\zeta|^{2 - {\rm Re}(\alpha)}} \leq \frac{A}{|\sin \theta|^{2 - {\rm Re}(\alpha)}}
 \]
 when $\cos \alpha < 0$. If ${\rm Re} (\alpha) - 2 >0$, then
 \[
 A |1+\zeta|^{{\rm Re}(\alpha) -2} \leq B {\rm e}^{|\zeta|}
 \]
 for an appropriate constant $B > 0$. Therefore, the function $\varphi(\zeta)$ is of exponential size at most~1 at~$\infty$
 along any direction $\theta \neq \pi$ from $0$ to $+\infty {\rm e}^{{\rm i} \theta}$. Moreover, the function $\varphi(\zeta)$ is
 continued analytically along any such a direction. Then the corresponding Laplace transform $(\mathcal{L}_{\theta} \varphi)(\tau)$
 is well defined and gives the 1-sum of the divergent power series $\hat{\varphi}(\tau)$ in such a direction. If we denote
 by $\varphi_{\theta}(\tau)$ this 1-sum, we have that
 \[
 \varphi_{\theta}(\tau)=2 (\alpha-1)
 \int_0^{+\infty {\rm e}^{{\rm i} \theta}}
 (1+\zeta)^{\alpha-2} {\rm e}^{-\frac{\zeta}{\tau}} {\rm d} \zeta
 \]
 for every $\theta \neq \pi$. From the above estimates, it follows that when ${\rm Re} (\alpha) -2 \leq 0$, the 1-sum
 $\varphi_{\theta}(\tau)$ is a holomorphic function in the open unlimited disc
 \begin{equation}
 \mathcal{D}_{\theta}=\biggl\{ \tau\in\mathbb{C} \mid {\rm Re} \biggl(\frac{{\rm e}^{{\rm i} \theta}}{\tau}\biggr) > 0 \biggr\} ,\label{d}
 \end{equation}
 whose opening is $\pi$.
 When ${\rm Re} (\alpha) - 2 > 0$, the 1-sum is a holomorphic function in the open bounded disc
 \begin{equation}
 \mathcal{D}_{\theta}(1)=\biggl\{ \tau\in\mathbb{C} \mid {\rm Re} \biggl(\frac{{\rm e}^{{\rm i} \theta}}{\tau}\biggr) > 1 \biggr\}\label{d1}
 \end{equation}
 with opening $ < \pi$.

 Using similar arguments, we find that for any direction $\theta \neq 0$ the Laplace transform
 \[
 \phi_{\theta}(\tau)=\int_0^{+\infty {\rm e}^{{\rm i} \theta}}
 (1-\zeta)^{-\alpha} {\rm e}^{-\frac{\zeta}{\tau}} {\rm d} \zeta
 \]
 defines the 1-sum of the divergent power series $\hat{\phi}(\tau)$ in such a direction. When ${\rm Re} (\alpha) \geq 0$,
 the 1-sum $\phi_{\theta}(\tau)$ is a holomorphic function in the disc $\mathcal{D}_{\theta}$ from~\eqref{d}. Otherwise, the
 1-sum is a~holomorphic function in the disc $\mathcal{D}_{\theta}(1)$ from~\eqref{d1}.
%This ends the proof.
 \end{proof}

\end{document}
